\documentclass[12pt,a4paper]{article}

\usepackage[utf8]{inputenc}
\usepackage[T1]{fontenc}
\usepackage[italian]{babel}
\usepackage{geometry}
\usepackage{longtable}
\usepackage{array}
\usepackage{booktabs}
\usepackage[hidelinks]{hyperref}

\geometry{margin=2.5cm}

\title{Gestionale Affitti\\Specifica dei requisiti}
\author{}
\date{}

\newcommand{\req}[1]{\textbf{#1}}

\begin{document}

\maketitle
\tableofcontents
\newpage

\section{Scopo}

\emph{Gestionale Affitti} è un progetto individuale per la gestione essenziale di
locazioni immobiliari.

La versione 1.0 è intenzionalmente limitata a due obiettivi principali:
\begin{itemize}
    \item registrare un contratto di locazione relativo all'intero immobile;
    \item registrare il pagamento mensile del canone relativo a un contratto.
\end{itemize}

L'unico attore che interagisce direttamente con il sistema nella versione 1.0 è il
\textbf{proprietario}. L'\textbf{inquilino} è un concetto del dominio ma non è un attore
del sistema.

La specifica dettagliata dei due casi d'uso e la relativa modellazione vengono rinviate
alla fase UML. I due casi d'uso core individuati sono:

\begin{itemize}
    \item \textbf{UC-01 -- Registrare un contratto di locazione};
    \item \textbf{UC-02 -- Registrare il pagamento di un canone}.
\end{itemize}

\section{Scope della versione 1.0}

Rientrano nello scope:
\begin{itemize}
    \item locazione dell'intero immobile;
    \item registrazione dei dati di immobile, proprietario e inquilino;
    \item contratti a canone concordato e a canone libero;
    \item clausole contrattuali predefinite per la tipologia selezionata;
    \item registrazione del pagamento completo del canone mensile;
    \item controlli necessari a evitare contratti e pagamenti incoerenti.
\end{itemize}

Sono fuori scope:
\begin{itemize}
    \item locazione parziale dell'immobile;
    \item rinnovo operativo dei contratti e avvisi di prossimità della scadenza;
    \item pagamenti parziali;
    \item esecuzione reale del pagamento e integrazione con servizi bancari;
    \item firma e registrazione fiscale del contratto;
    \item modifica o caricamento di clausole contrattuali da parte del proprietario;
    \item manutenzione, annunci e prenotazione degli immobili;
    \item autenticazione e autorizzazione degli utenti;
    \item utilizzo pubblico o multiutente del sistema.
\end{itemize}

La versione 1.0 assume pertanto un contesto dimostrativo e mono-utente, non destinato
all'esposizione pubblica o all'uso in produzione. L'utilizzatore opera nel ruolo di proprietario
e il sistema non deve distinguere dati appartenenti a utenti diversi; autenticazione e
autorizzazione restano fuori scope.

\section{Regole di dominio essenziali}

\begin{itemize}
    \item Ogni immobile può avere un nome assegnato dal proprietario. Il nome è un'etichetta
          leggibile, non deve essere univoco e non identifica l'immobile.
    \item L'immobile è riconosciuto tramite i dati catastali; non devono essere creati duplicati
          con gli stessi dati catastali.
    \item proprietario e inquilino sono riconosciuti tramite codice fiscale.
    \item Ogni contratto può avere un nome o una breve descrizione leggibile.
    \item La versione 1.0 supporta contratti a canone concordato (3 + 2) e a canone libero (4 + 4).
    \item Il proprietario inserisce la data iniziale \texttt{dal}; la data finale \texttt{al}
          è determinata automaticamente in base al contratto selezionato e rappresenta l'ultimo
          giorno compreso nel periodo iniziale.
    \item Il sistema gestisce esclusivamente il periodo iniziale del contratto; non gestisce
          operazioni o avvisi relativi al rinnovo successivo, che restano fuori scope nella versione 1.0.
    \item Il canone è mensile e il giorno di pagamento deve essere compreso tra 1 e 28.
    \item Un immobile può avere più contratti di locazione nel tempo, purché i relativi
          periodi di decorrenza non si sovrappongano.
    \item Non sono ammessi pagamenti parziali del canone di locazione; il relativo importo deriva dal canone mensile nel contratto.
    \item Un pagamento può essere registrato anche dopo la relativa scadenza.
\end{itemize}

\section{User stories}

\subsection{US-01 -- Registrazione del contratto}

\emph{Come proprietario, voglio registrare un contratto di locazione relativo a un immobile,
in modo da gestire il rapporto locativo con l'inquilino.}

\subsection{US-02 -- Registrazione del pagamento}

\emph{Come proprietario, voglio registrare il pagamento del canone mensile di un contratto,
in modo da tenere traccia delle mensilità pagate.}

\section{Requisiti funzionali}

\begin{longtable}{>{\raggedright\arraybackslash}p{1.8cm}
                  >{\raggedright\arraybackslash}p{13cm}}
\toprule
\textbf{ID} & \textbf{Requisito funzionale} \\
\midrule
\endfirsthead

\toprule
\textbf{ID} & \textbf{Requisito funzionale} \\
\midrule
\endhead

\req{RF-01} &
Il sistema deve consentire al proprietario di registrare un contratto mediante una
procedura guidata in sei step: immobile -> proprietario -> inquilino -> dati contrattuali ->
riepilogo -> conferma. I dati obbligatori di uno step devono essere validi prima di
proseguire e la procedura deve poter essere annullata prima della conferma definitiva.
Se alla registrazione è presente una bozza, il sistema carica i dati della stessa. \\

\req{RF-02} &
Al completamento del primo step il sistema deve creare una bozza del contratto e aggiornarla
al completamento di ogni step validato con successo. La bozza deve essere cancellata quando
il proprietario annulla la procedura oppure quando il contratto viene registrato definitivamente
con successo. \\

\req{RF-03} &
Durante la registrazione del contratto il sistema deve consentire di selezionare un immobile
già registrato tramite il nome assegnato oppure di inserirne uno nuovo. Gli immobili devono
essere riconosciuti tramite i dati catastali, mentre proprietario e inquilino devono essere
riconosciuti tramite codice fiscale. \\

\req{RF-04} &
Il sistema deve acquisire per il contratto un nome o una breve descrizione, la tipologia
contrattuale, la data iniziale, il canone mensile e il giorno di pagamento. Deve calcolare
automaticamente la data finale del periodo iniziale e associare le clausole predefinite
alla tipologia contrattuale selezionata. \\

\req{RF-05} &
La presenza di uno o più contratti già registrati per l'immobile non deve impedire l'avvio o
la prosecuzione della stesura di un nuovo contratto. Il sistema deve impedire la registrazione
definitiva soltanto quando il periodo di decorrenza del nuovo contratto si sovrappone a quello
di un altro contratto già registrato per lo stesso immobile. Lo stesso inquilino può quindi
stipulare, in periodi non sovrapposti, più contratti successivi sul medesimo immobile. \\

\req{RF-06} &
Prima della registrazione definitiva il sistema deve mostrare un riepilogo dei dati
contrattuali e consentire al proprietario di modificarli. Il contratto deve essere
registrato soltanto dopo una conferma esplicita.
Per registrazione si intende il salvataggio permanente dei dati e non la registrazione
ufficiale presso l'Agenzia delle Entrate. \\

\req{RF-07} &
Per registrare un pagamento il sistema deve consentire al proprietario di selezionare
l'immobile tramite il nome assegnato e mostrare, senza richiedere l'unicità del nome,
gli inquilini associati a contratti registrati per quell'immobile. Se l'elenco contiene
un solo inquilino, il sistema deve selezionarlo automaticamente. \\

\req{RF-08} &
Dopo l'individuazione di immobile e inquilino, il sistema deve considerare i contratti
registrati relativi alla coppia selezionata e individuare automaticamente la mensilità non
pagata cronologicamente più vecchia fino al mese corrente. Deve quindi caricare i dati
essenziali del contratto a cui tale mensilità appartiene, escludendo mensilità già pagate
e mensilità future. \\

\req{RF-09} &
L'importo del pagamento è determinato automaticamente dal canone mensile del
contratto. La versione 1.0 deve consentire soltanto il pagamento completo e deve permettere
la registrazione anche quando il pagamento è tardivo. \\

\req{RF-10} &
Il sistema deve offrire le azioni \emph{Pagato} e \emph{Annulla}. La registrazione del
pagamento deve avvenire soltanto dopo una conferma esplicita; in caso di annullamento
nessun pagamento deve essere registrato. \\

\req{RF-11} &
Nella home page il sistema deve mostrare un elenco degli immobili registrati con il nome
assegnato e il numero di contratti registrati. Per ciascun immobile deve inoltre evidenziare
lo stato complessivo dei pagamenti, indicando la regolarità dei pagamenti oppure il numero
di mensilità arretrate non ancora registrate come pagate. \\

\bottomrule
\end{longtable}

\section{Requisiti non funzionali}

\begin{longtable}{>{\raggedright\arraybackslash}p{1.5cm}
                  >{\raggedright\arraybackslash}p{10cm}
                  >{\raggedright\arraybackslash}p{2.2cm}}
\toprule
\textbf{ID} & \textbf{Requisito non funzionale} & \textbf{Categoria} \\
\midrule
\endfirsthead

\toprule
\textbf{ID} & \textbf{Requisito non funzionale} & \textbf{Categoria} \\
\midrule
\endhead

\req{RNF-01} &
I dati validati negli step completati della registrazione del contratto devono essere
recuperabili dopo un'interruzione della procedura o una temporanea indisponibilità del
servizio. Nel contesto mono-utente della versione 1.0, alla riapertura della procedura il
sistema deve poter recuperare l'eventuale bozza non conclusa senza richiedere un meccanismo
di identificazione tra utenti diversi. &
Affidabilità / recuperabilità \\

\req{RNF-02} &
In condizioni nominali di utilizzo, le operazioni interattive dei due casi d'uso che
richiedono elaborazione o comunicazione con il server devono completarsi entro
1,5 secondi. &
Prestazioni \\

\req{RNF-03} &
Il client deve applicare controlli di ammissibilità coerenti con il tipo di ciascun campo,
segnalando prima dell'invio caratteri o sequenze non consentiti. Tali controlli non sostituiscono
la validazione lato server: ogni input deve essere nuovamente validato e trattato in modo sicuro
prima di essere visualizzato o utilizzato nelle operazioni di persistenza, con particolare
attenzione a vulnerabilità di tipo XSS e injection. Se la validazione lato server rileva un input
riconducibile a un possibile tentativo di attacco, la richiesta deve essere rifiutata e deve essere
registrato un evento di sicurezza contenente i metadati tecnici disponibili della richiesta e il
valore rilevato in forma sanitizzata. &
Sicurezza \\

\req{RNF-04} &
Gli errori applicativi recuperabili devono essere intercettati e registrati nei log.
Un errore relativo a una singola operazione non deve lasciare dati parziali o incoerenti,
non deve produrre un falso esito positivo e non deve impedire le operazioni successive.
Il messaggio mostrato all'utente non deve esporre dettagli tecnici interni. &
Affidabilità / resilienza \\

\req{RNF-05} &
Durante le operazioni che richiedono attesa per elaborazione o comunicazione con il server,
il sistema deve rendere percepibile lo stato di elaborazione e, al termine, comunicare
esplicitamente l'esito positivo o negativo. &
Usabilità \\

\bottomrule
\end{longtable}

\section{Acceptance criteria}

\subsection{Acceptance criteria di US-01}

\paragraph{AC-01 -- Selezione o inserimento dell'immobile}
\textbf{Dato} che il proprietario avvia la registrazione di un contratto,
\textbf{quando} accede al primo step,
\textbf{allora} deve poter selezionare un immobile esistente tramite il nome assegnato
oppure inserire un nuovo immobile; un duplicato catastale deve essere rifiutato.

\paragraph{AC-02 -- immobile con contratti già registrati}
\textbf{Dato} un immobile per il quale esiste già uno o più contratti registrati,
\textbf{quando} il proprietario lo seleziona per predisporre un nuovo contratto,
\textbf{allora} il sistema deve consentire di proseguire con la stesura; la registrazione
definitiva verrà impedita soltanto se, una volta noto il nuovo periodo di decorrenza,
risulta una sovrapposizione con un contratto esistente.

\paragraph{AC-03 -- Identificazione delle persone}
\textbf{Dato} l'inserimento del codice fiscale del proprietario o dell'inquilino,
\textbf{quando} il soggetto è già registrato,
\textbf{allora} il sistema deve recuperarne i dati disponibili.

\paragraph{AC-04 -- Validazione progressiva}
\textbf{Dato} uno qualsiasi dei primi quattro step,
\textbf{quando} mancano dati obbligatori o sono presenti dati non validi,
\textbf{allora} il sistema deve segnalare gli errori e impedire il passaggio allo step
successivo finché non vengono corretti.

\paragraph{AC-05 -- Dati e durata del contratto}
\textbf{Dato} che il proprietario ha scelto la tipologia contrattuale e la data
\texttt{dal},
\textbf{quando} i dati contrattuali vengono validati,
\textbf{allora} il sistema deve calcolare automaticamente la data \texttt{al},
associare le clausole predefinite e accettare un giorno di pagamento compreso tra 1 e 28.

\paragraph{AC-06 -- Coerenza del contratto}
\textbf{Dato} il periodo del nuovo contratto,
\textbf{quando} tale periodo si sovrappone al periodo di un contratto già registrato per lo stesso
immobile,
\textbf{allora} il sistema deve impedire la registrazione definitiva del nuovo contratto.

\paragraph{AC-07 -- Riepilogo e conferma}
\textbf{Dato} che gli step precedenti sono validi,
\textbf{quando} il proprietario raggiunge il riepilogo,
\textbf{allora} deve poter verificare e modificare i dati; il contratto deve essere
registrato solo dopo la conferma esplicita.

\paragraph{AC-08 -- Annullamento}
\textbf{Dato} che il contratto non è ancora stato confermato,
\textbf{quando} il proprietario annulla la procedura,
\textbf{allora} nessun nuovo contratto deve essere registrato e l'eventuale bozza deve essere cancellata.

\paragraph{AC-15 -- Ciclo di vita della bozza}
\textbf{Dato} che il proprietario completa con successo il primo step della procedura,
\textbf{quando} passa a uno step successivo validato,
\textbf{allora} il sistema deve creare o aggiornare la bozza con i dati acquisiti;
\textbf{quando} la registrazione definitiva termina con successo, la bozza deve essere cancellata.

\subsection{Acceptance criteria di US-02}

\paragraph{AC-09 -- Individuazione dell'inquilino}
\textbf{Dato} che il proprietario avvia la registrazione di un pagamento,
\textbf{quando} seleziona un immobile,
\textbf{allora} il sistema deve mostrare gli inquilini associati a contratti registrati
per quell'immobile e auto-selezionare l'inquilino se è l'unico disponibile.

\paragraph{AC-10 -- Caricamento dei dati contrattuali}
\textbf{Dato} che sono stati individuati immobile, inquilino e mensilità da registrare,
\textbf{quando} il sistema prepara il pagamento,
\textbf{allora} deve mostrare almeno canone mensile, tipologia e periodo
\texttt{dal}--\texttt{al} del contratto a cui la mensilità appartiene.

\paragraph{AC-11 -- Prima mensilità non pagata}
\textbf{Dato} che per la coppia immobile--inquilino esistono uno o più contratti registrati,
\textbf{quando} esistono mensilità non ancora pagate fino al mese corrente,
\textbf{allora} il sistema deve selezionare automaticamente la mensilità non pagata
cronologicamente più vecchia tra i contratti considerati, senza proporre mensilità già pagate
o future.

\paragraph{AC-12 -- Importo del pagamento}
\textbf{Dato} che è stata individuata la mensilità da registrare,
\textbf{quando} il sistema prepara il pagamento,
\textbf{allora} l'importo deve coincidere automaticamente con il canone mensile previsto
dal contratto e non deve essere richiesto un importo manuale.

\paragraph{AC-13 -- Pagamento tardivo}
\textbf{Dato} che la mensilità individuata è già scaduta,
\textbf{quando} il proprietario registra il pagamento,
\textbf{allora} il sistema deve consentire comunque la registrazione.

\paragraph{AC-14 -- Conferma o annullamento}
\textbf{Dato} che il pagamento è pronto per la registrazione,
\textbf{quando} il proprietario seleziona \emph{Pagato},
\textbf{allora} il sistema deve chiedere conferma prima di registrarlo;
\textbf{quando} seleziona \emph{Annulla}, nessun pagamento deve essere registrato.

\paragraph{AC-16 -- Stato dei pagamenti in home}
\textbf{Dato} che nella home page sono presenti immobili registrati,
\textbf{quando} il sistema ne mostra l'elenco,
\textbf{allora} per ciascun immobile deve mostrare il nome assegnato, il numero di contratti
registrati e uno stato che evidenzi la regolarità dei pagamenti oppure il numero di mensilità
arretrate non ancora registrate come pagate.

\section{Criteri di verifica dei requisiti non funzionali}

Questi criteri non introducono nuove user stories, ma indicano come verificare i requisiti
trasversali della versione 1.0.

\begin{description}
    \item[CV-01 -- RNF-01] Dopo aver completato almeno uno step della registrazione del
    contratto, un'interruzione della procedura deve consentire, nel contesto mono-utente
    previsto dalla versione 1.0, di riaprire la procedura recuperando l'ultima bozza non
    conclusa e i dati già acquisiti.

    \item[CV-02 -- RNF-02] In condizioni nominali, ciascuna operazione interattiva di
    UC-01 e UC-02 che richiede il server deve completarsi entro 1,5 secondi.

    \item[CV-03 -- RNF-03] Un input contenente caratteri o sequenze non ammessi per il campo
    deve essere segnalato dal client. Lo stesso input, se inviato direttamente al server aggirando
    il controllo client, deve essere nuovamente validato e rifiutato o neutralizzato senza eseguire
    codice indesiderato né alterare le operazioni di persistenza. Se il server lo riconduce a un
    possibile tentativo di XSS o injection, deve produrre una voce di log con metadati tecnici della
    richiesta e valore rilevato in forma sanitizzata.

    \item[CV-04 -- RNF-04] Un errore applicativo recuperabile deve produrre una voce di
    log, non deve lasciare dati parziali, deve restituire un messaggio utente privo di
    dettagli tecnici e non deve impedire una successiva operazione valida.

    \item[CV-05 -- RNF-05] Durante un'operazione con attesa deve essere mostrato uno stato
    di elaborazione; al termine deve essere comunicato un esito positivo o negativo.
\end{description}

\section{Tracciabilità iniziale}

La tabella mantiene soltanto il collegamento necessario alla fase successiva.
La tracciabilità verso UML, componenti, codice e test verrà estesa progressivamente.

\begin{longtable}{>{\raggedright\arraybackslash}p{2cm}
                  >{\raggedright\arraybackslash}p{5cm}
                  >{\raggedright\arraybackslash}p{4cm}
                  >{\raggedright\arraybackslash}p{2.5cm}}
\toprule
\textbf{User story} & \textbf{Requisiti funzionali} & \textbf{Acceptance criteria} &
\textbf{Caso d'uso} \\
\midrule
\endfirsthead

\toprule
\textbf{User story} & \textbf{Requisiti funzionali} & \textbf{Acceptance criteria} &
\textbf{Caso d'uso} \\
\midrule
\endhead

US-01 & RF-01--RF-06 & AC-01--AC-08, AC-15 & UC-01 \\
US-02 & RF-07--RF-11 & AC-09--AC-14, AC-16 & UC-02 \\

\bottomrule
\end{longtable}

\section{Questioni aperte}

Prima di considerare definitivamente chiusa la specifica dei requisiti restano da chiarire:

\begin{itemize}
    \item la prima mensilità dovuta quando il contratto decorre nel corso del mese e la
          data \texttt{dal} è successiva al giorno mensile previsto per il pagamento;
    \item definire in modo riproducibile le condizioni nominali e l'ambiente di misura utilizzati
          per verificare il tempo massimo di risposta previsto da RNF-02;
    \item l'eventuale introduzione di un requisito di conformità normativa, da formulare
          soltanto dopo aver individuato e verificato le fonti normative italiane applicabili
          alle tipologie contrattuali e alle clausole predefinite.
\end{itemize}

\end{document}
